\chapter{Geometria biwymierna i modele minimalne powierzchni}
\label{ch:ga8-biwymierna}

\section{Dlaczego izomorfizm bywa zbyt silnym pojęciem?}
\label{sec:ga8-motivation}

Przy badaniu krzywych i powierzchni wielokrotnie zastępowaliśmy
punkt osobliwy nowymi punktami lub kierunkami. Rozdmuchanie
$\pi:\widetilde S\to S$ gładkiej powierzchni w punkcie
nie jest izomorfizmem: włókno nad środkiem ma postać $\PP^1$.
Mimo to poza środkiem $\pi$ jest izomorfizmem. Taka zmiana
jest \emph{biwymierna}. W tym rozdziale badamy jej najbardziej
konkretne konsekwencje: odwzorowania wymierne, przekształcenie
Cremony, rozdmuchania i kontrakcje krzywych.
W końcowej części rozdziału przejdziemy od tych
konkretnych przykładów do dwuwymiarowego programu
modeli minimalnych. Nie będziemy rozwijać jego
wyżej wymiarowej wersji.

Przez cały rozdział $k$ jest algebraicznie domkniętym ciałem
charakterystyki zero. Powierzchnie są gładkie, całkowite i
rzutowe nad $k$, chyba że w przykładzie zaznaczono inaczej.

\begin{definicja}[Odwzorowanie wymierne i biwymierne]
\label{def:ga8-rational}
Niech $X,Y$ będą całkowitymi rozmaitościami algebraicznymi
nad $k$. \emph{Odwzorowaniem wymiernym} $\varphi:X\dashrightarrow Y$
nazywamy klasę par $(U,\varphi_U)$, w której $U\subseteq X$
jest niepustym otwartym podzbiorem Zariskiego, a
$\varphi_U:U\to Y$ jest morfizmem; dwie pary są równoważne,
gdy ich morfizmy zgadzają się na pewnym niepustym
otwartym podzbiorze wspólnej dziedziny. Odwzorowanie jest
\emph{dominujące}, gdy jego obraz jest gęsty w $Y$.
Nazywamy je \emph{biwymiernym}, jeżeli istnieją niepuste
otwarte $U\subseteq X$, $V\subseteq Y$ takie, że
$\varphi_U:U\to V$ jest izomorfizmem.
\end{definicja}

\begin{uwaga}[Dlaczego wystarcza gęste otoczenie?]
Na rozmaitości całkowitej każdy niepusty otwarty
podzbiór jest gęsty. Odwzorowanie wymierne ignoruje
więc podzbiór, na którym nie potrafimy podać morfizmu,
ale pamięta działanie na części gęstej. Aby mówić o
\emph{izomorfizmie}, trzeba natomiast opisać wszystkie punkty.
Domena, na której dana reprezentacja jest morfizmem,
nie musi być największą domeną określoności.
\end{uwaga}

\begin{twierdzenie}[Kryterium przez ciała funkcji]
\label{tw:ga8-function-fields}
Niech $X,Y$ będą całkowitymi rozmaitościami nad $k$.
Każde dominujące odwzorowanie wymierne
$\varphi:X\dashrightarrow Y$ wyznacza zanurzenie ciał
$\varphi^*:k(Y)\hookrightarrow k(X)$.
Odwzorowanie $\varphi$ jest biwymierne wtedy i tylko
wtedy, gdy $\varphi^*$ jest izomorfizmem ciał nad $k$.
\end{twierdzenie}
\begin{proof}
Dominujący morfizm na niepustym otwartym podzbiorze
przenosi funkcję wymierną na jej złożenie z $\varphi$.
Jeśli funkcja znika po złożeniu, jest zerowa na gęstym
obrazie, zatem jest zerowa w $k(Y)$. To daje injektywność.
Jeśli $\varphi$ jest izomorfizmem otwartych podzbiorów,
wyznacza odwrotne izomorfizmy ich ciał funkcji, które
są dokładnie $k(X)$ i $k(Y)$.
W drugą stronę wybierzmy afiniczne niepuste otwarte
$U\subseteq X$ i $V\subseteq Y$. Generatory algebry
współrzędnych $k[V]$ jako elementy $k(X)$ są ilorazami
funkcji regularnych na $U$; na pewnym mniejszym
niepustym otwartym $U'$ ich mianowniki są niezerowe.
To definiuje morfizm wymierny $U'\to Y$.
W taki sam sposób izomorfizm odwrotny ciał funkcji
definiuje odwzorowanie wymierne $Y\dashrightarrow X$.
Złożenia tych odwzorowań mają na poziomie ciał funkcji
identyczność. Ponieważ $X$ i $Y$ są rozdzielone,
odwzorowania te zgadzają się z identycznościami na
odpowiednich niepustych otwartych podzbiorach.
Po ich zmniejszeniu uzyskujemy wzajemnie odwrotne
izomorfizmy otwartych podzbiorów.
\end{proof}

\begin{przyklad}[Rozdmuchanie nie zmienia ciała funkcji]
Dla $\pi:\operatorname{Bl}_0\A^2\to\A^2$ opisanej
w rozdziale~\ref{ch:ga5-rozdmuchania} mamy izomorfizm
po usunięciu punktu $0$ i dzielnika wyjątkowego.
Zatem $k(\operatorname{Bl}_0\A^2)=k(x,y)$.
Obecność całej prostej rzutowej nad jednym punktem
nie zmienia ciała funkcji wymiernych powierzchni.
\end{przyklad}

\section{Kwadratowe przekształcenie Cremony: obliczenie od początku}
\label{sec:ga8-cremona}

\begin{definicja}[Standardowe przekształcenie Cremony]
\label{def:ga8-cremona}
Odwzorowanie wymierne
\[
\sigma:\PP^2_k\dashrightarrow\PP^2_k,\qquad
[x:y:z]\longmapsto[yz:xz:xy]
\]
nazywamy standardowym kwadratowym przekształceniem Cremony.
\end{definicja}

\begin{stwierdzenie}[Punkty bazowe i odwrotność]
\label{st:ga8-cremona-basic}
Niech $k$ będzie ciałem, a
$\sigma:\PP^2_k\dashrightarrow\PP^2_k$ odwzorowaniem wymiernym
$[x:y:z]\mapsto[yz:xz:xy]$.
Odwzorowanie $\sigma$ jest nieokreślone dokładnie w punktach
$p_1=[1:0:0]$, $p_2=[0:1:0]$ i $p_3=[0:0:1]$.
Na otwartym zbiorze $T=\{xyz\ne0\}\subseteq\PP^2$
wyznacza izomorfizm $T\to T$, dla którego
$\sigma^{-1}=\sigma$. Na każdej prostej współrzędnych
poza punktami bazowymi jest stałe i przyjmuje
odpowiedni punkt współrzędnych.
\end{stwierdzenie}
\begin{proof}
Wszystkie trzy współrzędne $yz,xz,xy$ są zerowe
wtedy i tylko wtedy, gdy co najmniej dwie spośród
$x,y,z$ są równe zeru. Są to dokładnie trzy wskazane
punkty. Jeśli $xyz\ne0$, to
\[
\sigma^2([x:y:z])
=[(xz)(xy):(yz)(xy):(yz)(xz)]
=[x^2yz:xy^2z:xyz^2]=[x:y:z].
\]
Na prostej $x=0$ przy $yz\ne0$ obraz to $[1:0:0]$.
Dla pozostałych dwóch prostych otrzymujemy analogicznie
punkty $p_2,p_3$. Ostatni rachunek pokazuje również,
że punkty leżące na tych prostych (poza bazowymi)
nie mogą należeć do otwartego zbioru, na którym
$\sigma$ jest izomorfizmem na obraz otwarty.
\end{proof}

\begin{uwaga}[Po co trzy rozdmuchania?]
Nie da się rozszerzyć wzoru $[yz:xz:xy]$ do morfizmu
całej $\PP^2$ tylko przez \emph{nadanie} mu wartości
w trzech punktach bazowych: różne kierunki dojścia
do danego punktu prowadzą do różnych granicznych
wartości w przestrzeni rzutowej. Rozdmuchanie koduje
te kierunki jako punkty dzielnika wyjątkowego.
\end{uwaga}

\begin{stwierdzenie}[Rozwiązanie punktu bazowego w chartach]
\label{st:ga8-local-resolution}
Niech $p_1=[1:0:0]$. Na charcie $x\ne0$ z
współrzędnymi $(u=y/x,v=z/x)$ przekształcenie Cremony
ma postać $[uv:v:u]$. Po rozdmuchaniu $(u,v)=(0,0)$
na obu chartach jest morfizmem.
\end{stwierdzenie}
\begin{proof}
W pierwszym charcie rozdmuchania kładziemy $v=tu$.
Dla $u\ne0$ mamy
$[uv:v:u]=[tu^2:tu:u]=[tu:t:1]$.
Prawa strona ma sens także dla $u=0$ i opisuje
na dzielniku wyjątkowym punkty $[0:t:1]$.
W drugim charcie kładziemy $u=sv$. Wtedy
$[uv:v:u]=[sv^2:v:sv]=[sv:1:s]$,
co również rozszerza się na $v=0$.
Na przecięciu chartów (gdzie $s=1/t$)
oba wzory reprezentują ten sam punkt projektowy.
Po rozdmuchaniu $p_2$ i $p_3$ wykonujemy identyczny
rachunek z permutacją współrzędnych.
\end{proof}

\begin{center}
\begin{tikzpicture}[>=Stealth,scale=1.0]
 \coordinate (A) at (0,2.2); \coordinate (B) at (-2.45,-1.25);
 \coordinate (C) at (2.45,-1.25);
 \draw[thick] (A)--(B)--(C)--cycle;
 \foreach \P in {A,B,C}{\fill[blue!70!black] (\P) circle (2.6pt);}
 \node[above=5pt,font=\footnotesize] at (A) {$p_3$};
 \node[below left=3pt,font=\footnotesize] at (B) {$p_1$};
 \node[below right=3pt,font=\footnotesize] at (C) {$p_2$};
 \node[align=center,font=\footnotesize,text width=5cm] at (0,-2.15)
 {Trzy punkty bazowe i trzy\\ proste współrzędnych};
\end{tikzpicture}

{\small Rysunek 8.1. Schemat położenia punktów bazowych
standardowego przekształcenia Cremony. Krawędzie trójkąta
oznaczają trzy proste kontraktowane przez odwzorowanie.}\par
\end{center}

\section{Forma przecięć na rozdmuchaniu trzech punktów}
\label{sec:ga8-classes}

Niech $\pi:X\to\PP^2$ będzie rozdmuchaniem
trzech różnych niewspółliniowych punktów $p_1,p_2,p_3$.
Oznaczmy przez $H$ klasę przeciwobrazu prostej,
a przez $E_1,E_2,E_3$ klasy dzielników wyjątkowych.
W rozdziale~\ref{ch:ga6-powierzchnie} wyprowadziliśmy
\[
H^2=1,\qquad H\cdot E_i=0,\qquad
E_i^2=-1,\qquad E_i\cdot E_j=0\quad(i\ne j).
\]
Kanonika ma klasę $K_X=-3H+E_1+E_2+E_3$.

\begin{stwierdzenie}[Transformaty prostych przez pary punktów]
\label{st:ga8-hexagon}
Dla $1\le i<j\le3$ oznaczmy przez $L_{ij}$
transformatę ścisłą prostej łączącej $p_i,p_j$.
Wówczas $[L_{ij}]=H-E_i-E_j$, a
$L_{ij}^2=-1$ oraz $K_X\cdot L_{ij}=-1$.
Ponadto $L_{ij}\cdot E_i=L_{ij}\cdot E_j=1$,
$L_{ij}\cdot E_k=0$ dla $k\notin\{i,j\}$,
oraz pary różnych $L_{ij}$ są rozłączne.
\end{stwierdzenie}
\begin{proof}
Własność transformacji ścisłej z rozdziału 6 daje
$[L_{ij}]=H-E_i-E_j$, ponieważ prosta przechodzi
przez $p_i,p_j$ z krotnością jeden i przez trzeci punkt
nie przechodzi. Formę przecięć obliczamy algebraicznie:
$L_{ij}^2=1-1-1=-1$,
$K_X\cdot L_{ij}=-3+1+1=-1$.
Tak samo $L_{ij}\cdot E_i= -E_i^2=1$.
Dla różnych par, na przykład $(i,j)=(1,2)$ i
$(i,j)=(1,3)$, dostajemy
$(H-E_1-E_2)\cdot(H-E_1-E_3)=1+E_1^2=0$.
Obie krzywe są całkowite, różne i efektywne, więc
nieujemne lokalne krotności przecięć wymuszają
rozłączność.
\end{proof}

\begin{center}
\begin{tikzpicture}[>=Stealth,scale=0.95]
\foreach \i/\a/\txt in {
0/90/$E_1$,1/30/$L_{12}$,2/-30/$E_2$,
3/-90/$L_{23}$,4/-150/$E_3$,5/-210/$L_{13}$}{
 \coordinate (P\i) at (\a:2.05cm);
 \fill[blue!65!black] (P\i) circle (2.4pt);
 \node[font=\footnotesize] at (\a:2.72cm) {\txt};
}
\draw[thick] (P0)--(P1)--(P2)--(P3)--(P4)--(P5)--cycle;
\node[align=center,font=\footnotesize] at (0,0)
 {przecięcia\\równe $1$\\na krawędziach};
\end{tikzpicture}

{\small Rysunek 8.2. Graf przecięć sześciu krzywych
samoprzecięcia $-1$ na rozdmuchaniu trzech niewspółliniowych
punktów $\PP^2$. Krawędzie odpowiadają iloczynowi $1$,
a pozostałe pary różnych krzywych mają iloczyn $0$.}\par
\end{center}

\begin{stwierdzenie}[Układ liniowy rozwiązujący Cremonę]
\label{st:ga8-cremona-linear}
Niech $X$ będzie jak wyżej, a $H'$ klasą przeciwobrazu
prostej z docelowej $\PP^2$ przez morfizm
$\widetilde\sigma:X\to\PP^2$ uzyskany po rozwiązaniu
punktów bazowych. Wtedy
$H'=2H-E_1-E_2-E_3$, a $H'^2=1$.
Morfizm $\widetilde\sigma$ kontraktuje każdą z krzywych
$L_{12},L_{13},L_{23}$ i odwzorowuje każdą $E_i$
izomorficznie na prostą w docelowej płaszczyźnie.
\end{stwierdzenie}
\begin{proof}
Współrzędne $yz,xz,xy$ są formami jednorodnymi
stopnia dwa, każda znika w trzech punktach bazowych
co najmniej do rzędu jeden. Ich transformaty po
rozdmuchaniu wyznaczają układ $|2H-E_1-E_2-E_3|$.
Rachunek w chartach z dowodu
stwierdzenia~\ref{st:ga8-local-resolution} pokazuje,
że transformaty trzech przekrojów nie zerują się
jednocześnie, zatem układ jest bez punktów bazowych
i definiuje $\widetilde\sigma$.
Stąd $H'=2H-E_1-E_2-E_3$ i $H'^2=4-3=1$.
Dla przykładu $H'\cdot L_{12}=2-1-1=0$,
podobnie dla pozostałych. Możemy też sprawdzić
kontrakcję bezpośrednio: prosta $z=0$ poza punktami
bazowymi przechodzi przez $\sigma$ w punkt $p_3$.
Jej transformata ścisła jest więc kontraktowana.
Na $E_1$ w pierwszym charcie mapa ma postać
$t\mapsto[0:t:1]$, a w drugim $s\mapsto[0:1:s]$;
razem tworzą parametryzację docelowej prostej $x'=0$.
Analogicznie zachowują się $E_2,E_3$.
\end{proof}

\section{Krzywe wyjątkowe i odwracanie rozdmuchania}
\label{sec:ga8-contraction}

\begin{definicja}[Krzywa $(-1)$]
\label{def:ga8-minus-one}
Niech $S$ będzie gładką powierzchnią rzutową nad $k$.
Krzywą $(-1)$ nazywamy gładką całkowitą krzywą
$E\subseteq S$, która jest izomorficzna z $\PP^1_k$
i spełnia $E^2=-1$.
\end{definicja}

\begin{stwierdzenie}[Dzielnik wyjątkowy rozdmuchania]
\label{st:ga8-exceptional-curve}
Jeżeli $\pi:\widetilde S\to S$ jest rozdmuchaniem
gładkiej powierzchni rzutowej w punkcie $p$, to
$E=\pi^{-1}(p)$ jest krzywą $(-1)$.
Ponadto $K_{\widetilde S}\cdot E=-1$.
\end{stwierdzenie}
\begin{proof}
Z konstrukcji rozdmuchania $E\cong\PP^1$, a z
rozdziału~\ref{ch:ga6-powierzchnie} $E^2=-1$.
Korzystając ze wzoru $K_{\widetilde S}=\pi^*K_S+E$
oraz z równości $\pi^*K_S\cdot E=0$, otrzymujemy
$K_{\widetilde S}\cdot E=E^2=-1$.
\end{proof}

\begin{twierdzenie}[Kryterium Castelnuova — wynik zewnętrzny]
\label{tw:ga8-castelnuovo}
Niech $S$ będzie gładką rzutową powierzchnią nad
algebraicznie domkniętym ciałem charakterystyki zero,
a $E\subseteq S$ krzywą $(-1)$. Wtedy istnieje
gładka rzutowa powierzchnia $S'$ oraz morfizm
$\rho:S\to S'$, który jest izomorfizmem poza $E$,
kontraktuje $E$ do jednego gładkiego punktu i jest
izomorficzny z rozdmuchaniem $S'$ w tym punkcie.
\end{twierdzenie}

\begin{uwaga}[Dlaczego nie każda ujemna krzywa znika gładko?]
Krzywa o samoprzecięciu $-2$ nie może być dzielnikiem
wyjątkowym rozdmuchania \emph{jednego gładkiego punktu},
gdyż taki dzielnik ma samoprzecięcie $-1$.
Krzywe $(-2)$ mogą być kontraktowane do \emph{punktów osobliwych}.
Przykładem jest ujemna sekcja powierzchni Hirzebrucha
$\mathbb F_2$, którą można kontraktować do wierzchołka
stożka kwadratowego. Jest to inny rodzaj kontrakcji,
którego nie należy nazywać odwrotnością zwykłego
rozdmuchania gładkiego punktu.
\end{uwaga}

\begin{przyklad}[Dwie odwrotne operacje]
Zacznijmy od $\PP^2$ i rozdmuchajmy punkt $p$.
Otrzymujemy $\mathbb F_1$ z wyjątkową sekcją
$E\cong\PP^1$, $E^2=-1$. Kontrakcja tej sekcji
odtwarza $\PP^2$. Nie wszystkie krzywe $\mathbb F_1$
mają samoprzecięcie $-1$: włókno rzutowania
$\mathbb F_1\to\PP^1$ ma samoprzecięcie $0$.
Ten przykład odróżnia kontrakcję od zwykłego
„zmazywania dowolnej krzywej”.
\end{przyklad}

\section{Niezmienniki biwymierne i liczby zmieniające się po rozdmuchaniu}
\label{sec:ga8-invariants}

\begin{stwierdzenie}[Zmiana liczb przecięć przy jednym rozdmuchaniu]
\label{st:ga8-classes-blowup}
Niech $\pi:\widetilde S\to S$ będzie rozdmuchaniem
w punkcie gładkiej powierzchni rzutowej. Wtedy
$K_{\widetilde S}^2=K_S^2-1$, a w charakterystyce zero
$\chi(\mathcal O_{\widetilde S})=\chi(\mathcal O_S)$.
W szczególności $\deg c_2(T_{\widetilde S})=
\deg c_2(T_S)+1$.
\end{stwierdzenie}
\begin{proof}
Pierwsza równość wynika ze wzoru na kanonikę
$K_{\widetilde S}=\pi^*K_S+E$ oraz relacji
$\pi^*K_S\cdot E=0$, $E^2=-1$.
Niezmienniczość charakterystyki Eulera wynika z
$\pi_*\mathcal O_{\widetilde S}=\mathcal O_S$
i $R^j\pi_*\mathcal O_{\widetilde S}=0$ dla $j>0$;
jest to znany wynik o wyższych obrazach prostych przy
rozdmuchaniu gładkiego punktu, przyjęty w
przykładzie~\ref{prz:ga7-blowup-c2}.
Z wzoru Noether, udowodnionego w rozdziale 7
na podstawie Hirzebrucha--Riemanna--Rocha,
otrzymujemy $12\chi=K^2+\deg c_2$.
Po odjęciu wzorów dla obu powierzchni dostajemy
zmianę stopnia $c_2$ równą $+1$.
\end{proof}

\begin{przyklad}[Liczby po trzech rozdmuchaniach]
Dla $X=\operatorname{Bl}_{p_1,p_2,p_3}\PP^2$
mamy $K_X^2=9-3=6$ i $\deg c_2(T_X)=3+3=6$,
a $\chi(\mathcal O_X)=1$. Sześć krzywych $(-1)$
z poprzedniej sekcji tworzy heksagon przecięć;
liczby $K^2$ i $c_2$ nie pozwalają rozpoznać
ich indywidualnego położenia, ale dobrze
kontrolują kolejne rozdmuchania.
\end{przyklad}

\begin{twierdzenie}[Niezmienniczość $q$ i $p_g$ — wynik standardowy]
\label{tw:ga8-q-pg}
Jeżeli $S,S'$ są gładkimi spójnymi rzutowymi
powierzchniami nad $\C$ i są biwymierne, to
$h^1(S,\mathcal O_S)=h^1(S',\mathcal O_{S'})$ oraz
$h^0(S,K_S)=h^0(S',K_{S'})$.
\end{twierdzenie}
\begin{uwaga}[Skąd to wiadomo?]
Dla rozdmuchania pierwsza równość wynika z zanikania
wyższych obrazów prostych $\mathcal O$.
Druga zachodzi, ponieważ holomorficzne formy
najwyższego stopnia na gładkiej powierzchni
przechodzą jednoznacznie między powierzchnią i jej
rozdmuchaniem; można to sprawdzić w chartach, a
formalny dowód korzysta z przedłużania form.
Przejście od pojedynczego rozdmuchania do dowolnego
odwzorowania biwymiernego używa twierdzenia o
faktoryzacji morfizmów biwymiernych gładkich powierzchni
na rozdmuchania punktów. Nie przeprowadzamy tutaj
pełnego dowodu tego standardowego twierdzenia.
\end{uwaga}

\begin{przyklad}[Powierzchnie wymierne i K3]
Powierzchnie $\PP^2$ i $\mathbb F_1$ są biwymierne,
a mimo to nie są izomorficzne: $\mathbb F_1$ zawiera
krzywą o samoprzecięciu $-1$, natomiast $\PP^2$ nie.
Gładka kwartyka K3 w $\PP^3(\C)$ ma
$p_g=h^0(K)=1$, a $\PP^2$ ma $p_g=0$.
Zatem kwartyka K3 nie jest powierzchnią wymierną.
Wniosek nie wynika z samej różnicy $K^2$, ponieważ
ta liczba \emph{zmienia się} przy rozdmuchaniu.
\end{przyklad}

% =============================================================
% DWUWYMIAROWY PROGRAM MODELI MINIMALNYCH: UZUPELNIENIE
% =============================================================
\section{Równoważność numeryczna i stożek krzywych}
\label{sec:ga8-ne-cone}

\emph{Idea.} Odwzorowanie biwymierne może kontraktować krzywe,
ale nie dowolne. Na powierzchni ich istnienie rozpoznajemy,
porównując liczby przecięć z dzielnikiem kanonicznym.
Zamiast rozpatrywać wszystkie krzywe pojedynczo, grupujemy je
według zachowania względem dywizorów. Powstanie w ten sposób
stożek, którego szczególne kierunki wskazują możliwe kontrakcje.

\begin{definicja}[Równoważność numeryczna]
\label{def:ga8-numerical}
Niech $S$ będzie gładką rzutową powierzchnią nad algebraicznie
domkniętym ciałem $k$ charakterystyki zero. Dwa dywizory
$D,D'$ na $S$ są \emph{numerycznie równoważne}, co zapisujemy
$D\equiv D'$, gdy $(D-D')\cdot C=0$ dla każdej całkowitej
krzywej $C\subseteq S$. Dwa cykle jednowymiarowe $Z,Z'$
są numerycznie równoważne, gdy $D\cdot(Z-Z')=0$ dla
każdego dywizora $D$ na $S$. Oznaczamy przez
$N^1(S)_{\R}$ i $N_1(S)_{\R}$ odpowiednie przestrzenie
ilorazowe po rozciągnięciu skalarów do $\R$.
Ich wspólny wymiar $\rho(S)$ nazywamy \emph{liczbą Picarda}.
Przez liczbę przecięcia otrzymujemy parowanie
$N^1(S)_{\R}\times N_1(S)_{\R}\to\R$.
\end{definicja}

\begin{uwaga}[Co jest twierdzeniem zewnętrznym?]
Skończoność wymiaru obu przestrzeni jest konsekwencją
twierdzenia Nérona--Severiego. Z definicji ilorazów
wynika niezdegenerowanie parowania; utożsamienie obu
wymiarów w kontekście powierzchni można uzyskać także
z symetrii i niezdegenerowania formy przecięć.
Nie utożsamiamy równoważności numerycznej z liniową:
na krzywej eliptycznej dwa różne punkty mają ten sam
stopień, a więc są numerycznie równoważne jako dywizory,
lecz na ogół nie są liniowo równoważne.
\end{uwaga}

\begin{definicja}[Stożek krzywych i dywizor nef]
\label{def:ga8-nef}
Niech $S$ będzie jak wyżej. Przez $\operatorname{NE}(S)$
oznaczamy stożek w $N_1(S)_{\R}$ składający się ze
skończonych sum $\sum_i a_i[C_i]$, gdzie $a_i\ge0$,
a $C_i\subseteq S$ są całkowitymi krzywymi.
Jego domknięcie oznaczamy $\overline{\operatorname{NE}}(S)$.
Dywizor rzeczywisty $D$ na $S$ nazywamy \emph{nef}
(ang. \emph{numerically effective}), gdy $D\cdot C\ge0$
dla każdej całkowitej krzywej $C\subseteq S$.
\emph{Promieniem ekstremalnym} stożka wypukłego
$\Gamma\subseteq N_1(S)_{\R}$ nazywamy promień
$R=\R_{\ge0}v\subseteq\Gamma$, $v\ne0$, taki że
z $u,w\in\Gamma$ oraz $u+w\in R$ wynika $u,w\in R$.
\end{definicja}

\begin{przyklad}[Dlaczego nie wystarcza sprawdzić $D^2$?]
Na $S=\operatorname{Bl}_p\PP^2$ niech $H$ będzie
przeciwobrazem klasy prostej, a $E$ dzielnikiem wyjątkowym.
Wtedy $H^2=1$, $H\cdot E=0$, $E^2=-1$.
Dywizor $D=H+E$ ma $D^2=0$, ale $D\cdot E=-1$,
więc nie jest nef. Natomiast $H$ jest nef, ponieważ
$H=\pi^*\mathcal O_{\PP^2}(1)$ i dla każdej krzywej
$C$ mamy $H\cdot C\ge0$; dla $C=E$ liczba ta wynosi zero.
To pokazuje, dlaczego nieujemność samoprzecięcia
dywizora nie jest kryterium nefności.
\end{przyklad}

\begin{stwierdzenie}[Stożek krzywych na $\mathbb F_n$]
\label{st:ga8-fn-cone}
Niech $n\ge0$, a $\mathbb F_n$ będzie powierzchnią
Hirzebrucha z rozdziału~\ref{ch:ga6-powierzchnie}.
Oznaczmy przez $s$ jej standardową sekcję, a przez $f$
włókno projekcji do $\PP^1$, przy czym
$s^2=-n$, $s\cdot f=1$, $f^2=0$.
Wówczas $\overline{\operatorname{NE}}(\mathbb F_n)$
jest stożkiem generowanym przez $[s]$ i $[f]$.
Ponadto $K_{\mathbb F_n}=-2s-(n+2)f$,
a więc $K\cdot s=n-2$ oraz $K\cdot f=-2$.
\end{stwierdzenie}
\begin{proof}
Używamy znanego opisu $\operatorname{Pic}(\mathbb F_n)
\cong\Z[s]\oplus\Z[f]$ (wynik z konstrukcji
wiązki rzutowej). Niech $C\subseteq\mathbb F_n$
będzie całkowitą krzywą różną od $s$ i od włókna.
W klasie numerycznej piszemy $[C]=a[s]+b[f]$.
Dla przecięcia z ogólnym włóknem otrzymujemy
$a=C\cdot f\ge0$. Ponieważ $C$ i $s$ nie mają
wspólnej składowej, $C\cdot s=b-na\ge0$,
a zatem $b\ge na\ge0$. Dla $n=0$ ten sam argument
daje $a,b\ge0$. Włókno i sekcja są skutecznymi
krzywymi, więc każdy generator należy do stożka.
Stąd całość stożka jest dokładnie ich nieujemną
powłoką, która jest już domknięta.
Wzór kanoniczny wynika z adjunkcji dla włókna
($K\cdot f=-2$) i dla sekcji $s\simeq\PP^1$:
$(K+s)\cdot s=-2$. Gdy $K=as+bf$, pierwsza
równość daje $a=-2$, a druga $(-2s+bf+s)\cdot s
=n+b=-2$, czyli $b=-n-2$.
\end{proof}

\begin{center}
\begin{tikzpicture}[>=Stealth,scale=1]
  \draw[->,gray!75] (-0.4,0)--(5.4,0) node[right,font=\footnotesize] {$[s]$};
  \draw[->,gray!75] (0,-0.3)--(0,3.7) node[above,font=\footnotesize] {$[f]$};
  \fill[blue!12] (0,0)--(4.8,0)--(4.8,3.15)--(0,3.15)--cycle;
  \draw[blue!70!black,very thick] (0,0)--(4.8,0);
  \draw[blue!70!black,very thick] (0,0)--(0,3.2);
  \node[font=\footnotesize,align=center,fill=white,inner sep=3pt]
  at (2.65,2.25) {$\overline{\operatorname{NE}}(\mathbb F_n)$\\
  dodatnie kombinacje $[s]$ i $[f]$};
  \fill (0,0) circle (1.5pt);
\end{tikzpicture}

{\small Rysunek 8.3. Schemat stożka krzywych powierzchni
Hirzebrucha w bazie $([s],[f])$. Rysunek przedstawia
wycinek nieograniczonego stożka; długości wektorów
nie reprezentują liczb przecięć.}\par
\end{center}

\section{Krzywe z ujemnym przecięciem kanonicznym}
\label{sec:ga8-negative-curves}

Poniższy lemat jest jednym z najważniejszych elementarnych
ogniw programu modeli minimalnych dla powierzchni.
Pokazuje, że krzywe o ujemnym samoprzecięciu, które
jednocześnie mają ujemne przecięcie z $K_S$, są
ściśle określonej postaci.

\begin{twierdzenie}[Elementarne rozpoznawanie krzywej $(-1)$]
\label{tw:ga8-negative-implies-minusone}
Niech $S$ będzie gładką rzutową powierzchnią nad
algebraicznie domkniętym ciałem $k$, a $C\subseteq S$
całkowitą krzywą. Jeżeli $C^2<0$ i $K_S\cdot C<0$,
to $C$ jest gładką krzywą wymierną oraz
$C^2=K_S\cdot C=-1$. W szczególności
$C\simeq\PP^1_k$.
\end{twierdzenie}
\begin{proof}
Dla całkowitej, niekoniecznie gładkiej krzywej
na gładkiej powierzchni wzór adjunkcji arytmetycznej
przyjmuje postać
$2p_a(C)-2=C^2+K_S\cdot C$, gdzie
$p_a(C)=1-\chi(\mathcal O_C)$.
Krzywa jest całkowita i rzutowa, więc
$H^0(C,\mathcal O_C)=k$; dlatego
$p_a(C)=h^1(C,\mathcal O_C)\ge0$.
Obie liczby $C^2$ i $K_S\cdot C$ są ujemnymi
liczbami całkowitymi, zatem ich suma jest nie
większa niż $-2$. Wzór adjunkcji daje równocześnie
$C^2+K_S\cdot C=2p_a(C)-2\ge-2$.
Musi więc zachodzić równość $-2$ i każda z obu
liczb równa się $-1$. Ponadto $p_a(C)=0$.
Dla normalizacji $\nu:\widetilde C\to C$ zachodzi
$p_a(C)=g(\widetilde C)+\sum_{p\in C}\delta_p$,
gdzie $\delta_p\ge0$ mierzy defekt normalności.
Oba składniki znikają. Zatem $C$ jest normalna,
a krzywa normalna jest gładka; jej rodzaj wynosi zero.
Nad algebraicznie domkniętym ciałem otrzymujemy
$C\cong\PP^1$.
\end{proof}

\begin{przyklad}[Dlaczego nie wystarcza $K_S\cdot C<0$?]
Na $\PP^2$ każda prosta $L$ ma $L^2=1$ i
$K_{\PP^2}\cdot L=-3$, lecz nie jest wyjątkową
krzywą rozdmuchania. Na $\mathbb F_n$ włókno $f$
spełnia $f^2=0$ i $K\cdot f=-2$; również nie
można go kontraktować do gładkiego punktu.
Dopiero dodatkowe założenie $C^2<0$ pozwala
zastosować poprzednie twierdzenie.
\end{przyklad}

\begin{twierdzenie}[Kryterium kontrakcji Castelnuova — wynik zewnętrzny]
\label{tw:ga8-castelnuovo-mmp}
Niech $S$ będzie gładką rzutową powierzchnią nad
algebraicznie domkniętym ciałem charakterystyki zero.
Jeżeli $E\subset S$ jest krzywą izomorficzną z $\PP^1$
i $E^2=-1$, to istnieje gładka rzutowa powierzchnia
$S'$ i morfizm $\pi:S\to S'$, będący rozdmuchaniem
$S'$ w jednym punkcie, dla którego $E$ jest dokładnie
dzielnikiem wyjątkowym.
\end{twierdzenie}

\begin{uwaga}[Ujemność samoprzecięcia nie daje dowolnej kontrakcji]
Krzywa o samoprzecięciu $-2$ może być kontraktowalna
do \emph{punktu osobliwego}, ale nie spełnia kryterium
Castelnuova kontrakcji do gładkiego punktu.
Na $\mathbb F_2$ sekcja $s$ ma $s^2=-2$,
a $K\cdot s=0$. Nie należy więc interpretować
wszystkich krzywych ujemnych jako kolejnych kroków
programu gładkich modeli minimalnych.
\end{uwaga}

\section{Modele minimalne i wymiar Kodairy}
\label{sec:ga8-kodaira-minimal}

\begin{definicja}[Minimalność gładkiej powierzchni]
\label{def:ga8-minimal}
Gładką rzutową powierzchnię $S$ nazywamy
\emph{minimalną w sensie gładkich powierzchni},
gdy nie zawiera żadnej krzywej $E\simeq\PP^1$
o samoprzecięciu $E^2=-1$. Jeśli $K_S$ jest nef,
to $S$ nazywamy \emph{modelem minimalnym w sensie
kanonicznym}. Te określenia nie są bez dodatkowych
hipotez równoważne. Przymiotnik \emph{względnie minimalny}
rezerwujemy natomiast dla powierzchni z ustalonym morfizmem
do bazy, kiedy minimalność bada się względem tego morfizmu.
\end{definicja}

\begin{stwierdzenie}[Istnienie powierzchni minimalnej w sensie gładkim]
\label{st:ga8-terminate-minusone}
Niech $S$ będzie gładką rzutową powierzchnią nad
algebraicznie domkniętym ciałem charakterystyki zero.
Istnieje skończony ciąg kontrakcji krzywych $(-1)$
prowadzący do gładkiej powierzchni minimalnej w sensie gładkim.
\end{stwierdzenie}
\begin{proof}
Jeżeli $S$ nie zawiera krzywych $(-1)$, to koniec.
W przeciwnym wypadku kryterium Castelnuova dostarcza
morfizmu $S\to S_1$ kontraktującego taką krzywą.
Rozdmuchanie punktu zwiększa liczbę Picarda o jeden:
$\rho(S)=\rho(S_1)+1$; pokazaliśmy już rozkład
klas po rozdmuchaniu w rozdziale 6. Powierzchnia
$S_1$ znowu jest gładka i rzutowa. Powtarzając
konstrukcję, za każdym razem zmniejszamy dodatnią
liczbę całkowitą $\rho$. Nie da się więc wykonać
nieskończenie wielu kroków. Ostatnia powierzchnia
nie ma krzywych $(-1)$.
\end{proof}

\begin{przyklad}[Minimalność nie oznacza jednoznaczności]
Powierzchnia $\PP^2$ jest minimalna w sensie gładkim,
ponieważ każda niezerowa efektywna krzywa ma
klasę $dH$ z $d\ge1$, więc jej samoprzecięcie
to $d^2>0$. Powierzchnia
$\PP^1\times\PP^1$ również nie zawiera krzywych
$(-1)$: dla całkowitej krzywej klasy $aF_1+bF_2$
mamy $a,b\ge0$ i $C^2=2ab\ge0$.
Jednocześnie obie powierzchnie są wymierne, a więc
biwymierne, chociaż nie są izomorficzne:
ich liczby Picarda wynoszą odpowiednio $1$ i $2$.
To pokazuje, że bez dodatkowych hipotez
minimalny w sensie gładkim przedstawiciel klasy
biwymierności nie jest jedyny.
\end{przyklad}

\begin{definicja}[Wielowymiarowe układy kanoniczne i wymiar Kodairy]
\label{def:ga8-kodaira}
Niech $S$ będzie gładką spójną rzutową powierzchnią
nad $\C$. Dla $m\ge1$ definiujemy \emph{$m$-ty
plurigenus} przez $P_m(S)=h^0(S,\omega_S^{\otimes m})$.
Jeżeli $P_m(S)=0$ dla wszystkich $m\ge1$, przyjmujemy
$\kappa(S)=-\infty$. W przeciwnym razie
\emph{wymiarem Kodairy} nazywamy
$\kappa(S)=\max_{m:P_m(S)>0}\dim \varphi_{|mK_S|}(S)$,
gdzie $\varphi_{|mK_S|}$ jest odwzorowaniem wymiernym
danym pełnym układem liniowym $|mK_S|$; gdy układ
ma jeden niezależny przekrój, jego obraz jest punktem.
Dla powierzchni $\kappa(S)\in\{-\infty,0,1,2\}$.
\end{definicja}

\begin{stwierdzenie}[Niezmienniczość plurigenera przy rozdmuchaniu]
\label{st:ga8-plurigenera-blowup}
Niech $\pi:\widetilde S\to S$ będzie rozdmuchaniem
punktu gładkiej rzutowej powierzchni $S$ nad $\C$.
Dla każdego $m\ge1$ zachodzi
$H^0(\widetilde S,\omega_{\widetilde S}^{\otimes m})
\cong H^0(S,\omega_S^{\otimes m})$.
W konsekwencji wszystkie $P_m$ i wymiar Kodairy
są niezmiennikami biwymiernymi gładkich
rzutowych powierzchni.
\end{stwierdzenie}
\begin{proof}
Wiemy, że $K_{\widetilde S}=\pi^*K_S+E$,
gdzie $E$ jest dzielnikiem wyjątkowym. Dla $m\ge1$
wzór projekcji daje
$\pi_*\mathcal O_{\widetilde S}(mK_{\widetilde S})
=\mathcal O_S(mK_S)\otimes\pi_*\mathcal O_{\widetilde S}(mE)$.
Pokażemy, że $\pi_*\mathcal O_{\widetilde S}(mE)
=\mathcal O_S$. Nad $S\setminus\{p\}$ obie wiązki
są identyczne. Sekcja $\mathcal O(mE)$ nad
$\pi^{-1}(U)$ może być utożsamiona z funkcją wymierną,
która jest regularna nad $U\setminus\{p\}$:
warunek na dywizor nie dopuszcza biegunów poza $E$.
Ponieważ gładka powierzchnia $S$ jest normalna,
regularne funkcje przedłużają się przez podzbiór
kowymiaru dwa. Tak otrzymujemy sekcję $\mathcal O_U$.
Odwrotnie każda funkcja regularna na $U$ daje
sekcję $\mathcal O(mE)$. Wzór projekcji i wzięcie
sekcji globalnych kończą dowód równości plurigenusów.
Co więcej, po pomnożeniu pullbacku przekroju $mK_S$
przez kanoniczny przekrój $\mathcal O(mE)$ otrzymujemy
\emph{wszystkie} przekroje $mK_{\widetilde S}$. Dlatego
jeśli $|mK_S|$ jest niepusty, układ $|mK_{\widetilde S}|$
powstaje przez pomnożenie pullbacków wszystkich przekrojów
$|mK_S|$ przez ten sam przekrój $\mathcal O(mE)$.
Na otwartym podzbiorze, na którym $\pi$ jest izomorfizmem,
oba odwzorowania plurikanoniczne zgadzają się po utożsamieniu
dziedzin i mają zatem ten sam wymiar obrazu.
To osobno uzasadnia niezmienniczość $\kappa$ z definicji
przez wymiar obrazów, a nie tylko równość liczb $P_m$.
Przejście od rozdmuchań do dowolnej biwymierności
korzysta z twierdzenia o faktoryzacji odwzorowań
biwymiernych gładkich rzutowych powierzchni,
które przyjmujemy jako wynik zewnętrzny.
\end{proof}

\begin{przyklad}[Wszystkie cztery wartości wymiaru Kodairy]
\label{prz:ga8-kodaira-four}
Rozpatrzmy gładkie krzywe projektowe $C_g,C_h$
rodzajów $g,h\ge2$ i krzywą eliptyczną $E$.
\begin{enumerate}[label=(\roman*)]
\item $\kappa(\PP^2)=-\infty$, ponieważ
$K_{\PP^2}=-3H$ i $H^0(\PP^2,\mathcal O(-3m))=0$
dla każdego $m\ge1$.
\item Jeżeli $S$ jest K3, to $K_S\sim0$ i
$P_m(S)=1$ dla wszystkich $m\ge1$, zatem $\kappa(S)=0$.
\item Dla $S=C_g\times E$ mamy
$\omega_S\simeq\mathrm{pr}_1^*\omega_{C_g}$,
a więc $P_m(S)=h^0(C_g,\omega_{C_g}^{\otimes m})$.
Dla $m\ge2$ wzór Riemanna--Rocha na krzywej
daje $(2m-1)(g-1)$, zaś odwzorowania
plurikanoniczne mają obraz będący krzywą.
Stąd $\kappa(S)=1$.
\item Dla $S=C_g\times C_h$ mamy
$P_m(S)=(2m-1)^2(g-1)(h-1)$ dla $m\ge2$.
Wymiar obrazu odpowiednich odwzorowań wynosi dwa,
więc $\kappa(S)=2$.
\end{enumerate}
W dwóch przykładach z iloczynami użyliśmy równości
$H^0(C\times D,L\boxtimes M)
\cong H^0(C,L)\otimes H^0(D,M)$ oraz faktu,
że dostatecznie wysoka wielokrotność dywizora
kanonicznego na krzywej rodzaju co najmniej dwa
daje odwzorowanie o jednowymiarowym obrazie.
\end{przyklad}

\section{Dwuwymiarowy program modeli minimalnych}
\label{sec:ga8-surface-mmp}

Przechodzimy od procedury kontraktowania wszystkich
krzywych $(-1)$ do kryterium, które samo wskazuje,
\emph{czy} i \emph{jakiego typu} należy wykonać następny krok.
To jest zasadnicza idea programu modeli minimalnych.

\begin{twierdzenie}[Stożek i kontrakcje dla gładkich powierzchni — wynik zewnętrzny]
\label{tw:ga8-surface-contraction}
Niech $S$ będzie gładką rzutową powierzchnią nad
algebraicznie domkniętym ciałem charakterystyki zero.
Część stożka $\overline{\operatorname{NE}}(S)$, na której
$K_S$ jest ujemny, jest kontrolowana przez przeliczalną
rodzinę promieni ekstremalnych generowanych przez klasy
krzywych wymiernych (twierdzenie o stożku).
Dla każdego $K_S$-ujemnego promienia ekstremalnego $R$
istnieje kontrakcja $f:S\to Y$ o spójnych włóknach,
kontraktująca dokładnie te krzywe $C$, dla których
$[C]\in R$. Jeśli $Y$ ma wymiar dwa, $f$ jest
kontrakcją krzywej $(-1)$ do gładkiego punktu.
Jeżeli $\dim Y<2$, to $f$ jest \emph{fibracją Moriego}:
ma dodatni względny wymiar, względną liczbę Picarda
równą jeden i $-K_S$ jest względnie obszerny.
\end{twierdzenie}

\begin{uwaga}[Zakres twierdzenia]
Pełny dowód wykorzystuje twierdzenie o stożku i twierdzenie
o kontrakcjach. Nie wyprowadzamy ich z samego twierdzenia
Castelnuova. W przypadku $\dim Y=1$ ogólne włókno
jest izomorficzne z $\PP^1$; dla $\dim Y=0$ dywizor
$-K_S$ jest obszerny, a $\rho(S)=1$.
To są trzy typy kroków: kontrakcja biwymierna,
fibracja nad krzywą lub kontrakcja do punktu.
\end{uwaga}

\begin{twierdzenie}[Program modeli minimalnych dla gładkich powierzchni]
\label{tw:ga8-surface-mmp}
Dla każdej gładkiej rzutowej powierzchni $S$
nad algebraicznie domkniętym ciałem charakterystyki zero
istnieje skończony ciąg kontrakcji $K$-ujemnych
promieni ekstremalnych. Kończy się on w jednej
i tylko jednej z następujących sytuacji:
\begin{enumerate}[label=(\alph*)]
\item otrzymujemy gładką rzutową powierzchnię $S_{\min}$,
dla której $K_{S_{\min}}$ jest nef;
\item otrzymujemy fibrację Moriego $f:S_r\to Y$,
gdzie $\dim Y\in\{0,1\}$.
\end{enumerate}
Jeśli $\kappa(S)\ge0$, zachodzi przypadek (a).
W przypadku (b) powierzchnia ma $\kappa=-\infty$.
\end{twierdzenie}
\begin{proof}
Korzystamy z zewnętrznego twierdzenia
\ref{tw:ga8-surface-contraction}. Gdy $K$ jest nef,
procedura się kończy. W przeciwnym razie istnieje
$K$-ujemny promień ekstremalny i jego kontrakcja.
Jeśli jest to fibracja Moriego, otrzymujemy (b)
i kończymy. Jeśli kontrakcja jest biwymierna,
odwraca rozdmuchanie punktu, więc nowe $S_1$
jest gładkie i $\rho(S_1)=\rho(S)-1$.
Ograniczoność dodatniej liczby $\rho$ daje skończoność
liczby kroków biwymiernych. Wynikiem jest więc
(a) lub (b). Jeśli zachodzi (b), to ogólne włókno
nad krzywą jest $\PP^1$ albo $-K$ jest obszerny
w przypadku bazy punktowej; w obu przypadkach
nie ma niezerowych przekrojów $mK$ dla $m>0$.
W pierwszym przypadku widać to po restrykcji
do ogólnego włókna, gdzie
$\deg(mK|_{\PP^1})=-2m<0$.
W drugim niezerowy efektywny $mK$ miałby
ujemne przecięcie z obszernym $-K$, co jest
niemożliwe. Niezmienniczość plurigenusów przy
krokach biwymiernych daje $\kappa(S)=-\infty$.
Tym samym gdy $\kappa(S)\ge0$, przypadek (b)
jest wykluczony.
\end{proof}

\begin{center}
\begin{tikzpicture}[>=Stealth,
 every node/.style={font=\footnotesize,align=center},
 nodebox/.style={draw,rounded corners,inner sep=4pt}]
\node[nodebox,text width=37mm] (start) at (0,3.1)
 {gładka powierzchnia $S$};
\node[nodebox,text width=37mm] (decision) at (0,1.65)
 {czy $K_S$ jest nef?};
\node[nodebox,text width=36mm] (nef) at (-3.2,-0.2)
 {tak: model kanonicznie minimalny};
\node[nodebox,text width=41mm] (ray) at (3.25,-0.2)
 {nie: kontrakcja promienia $K$-ujemnego};
\node[nodebox,text width=34mm] (blow) at (0.65,-2.5)
 {biwymierna: $\rho$ maleje; powtórz};
\node[nodebox,text width=34mm] (fibre) at (5.1,-2.5)
 {fibracja Moriego: koniec};
\draw[->] (start)--(decision);
\draw[->] (decision)--node[above left,font=\scriptsize] {tak} (nef);
\draw[->] (decision)--node[above right,font=\scriptsize] {nie} (ray);
\draw[->] (ray)--(blow);
\draw[->] (ray)--(fibre);
\end{tikzpicture}

{\small Rysunek 8.4. Schemat dwuwymiarowego programu modeli
minimalnych. W przypadku kontrakcji biwymiernej
wracamy do testu nefności na nowej powierzchni.}\par
\end{center}

\begin{przyklad}[Trzy możliwe zachowania promienia ekstremalnego]
\label{prz:ga8-extremal-examples}
\begin{enumerate}[label=(\roman*)]
\item Na $\operatorname{Bl}_p\PP^2=\mathbb F_1$
krzywa wyjątkowa $s$ spełnia $s^2=-1$,
$K\cdot s=-1$ i jej kontrakcja prowadzi do $\PP^2$.
\item Na $\mathbb F_2$ włókno $f$ spełnia
$f^2=0$, $K\cdot f=-2$. Kontrakcja odpowiada
projekcji $\mathbb F_2\to\PP^1$; nie jest
biwymierna i nie usuwa sekcji $s$ o $s^2=-2$.
\item Na $\PP^2$ promień klasy prostej $H$
ma $H^2=1$, $K\cdot H=-3$.
Kontrakcja prowadzi do punktu, a nie do innej
powierzchni. Jest to przypadek Fano o $\rho=1$.
\end{enumerate}
W szczególności nie należy utożsamiać każdego
$K$-ujemnego promienia ze znikającą krzywą $(-1)$.
\end{przyklad}

\begin{przyklad}[Dwie różne kontrakcje na jednej powierzchni]
\label{prz:ga8-f1-two-rays}
Niech $S=\mathbb F_1\simeq\operatorname{Bl}_p\PP^2$, a $s$
niech oznacza krzywą wyjątkową oraz $f$ włókno
naturalnego rzutowania $S\to\PP^1$. Formuły
ze stwierdzenia~\ref{st:ga8-fn-cone} dla $n=1$ dają
$K_S=-2s-3f$, $K_S\cdot s=-1$ i $K_S\cdot f=-2$.
Oba promienie stożka krzywych są więc $K$-ujemne.
Kontrakcja promienia $[s]$ jest biwymierna i daje
$S\to\PP^2$, czyli odwraca rozdmuchanie punktu.
Kontrakcja promienia $[f]$ ma włókna $\PP^1$
i daje morfizm $S\to\PP^1$; jest fibracją Moriego,
a nie morfizmem biwymiernym. Liczba Picarda włókien
względem bazy wynosi jeden, a $-K_S\cdot f=2>0$.
Dlatego wybór promienia ekstremalnego naprawdę wpływa
na typ kroku i wynik programu, gdy $\kappa=-\infty$.
\end{przyklad}

\begin{twierdzenie}[Jednoznaczność dla nieujemnego wymiaru Kodairy — wynik zewnętrzny]
\label{tw:ga8-minimal-unique}
Niech $S,S'$ będą gładkimi rzutowymi powierzchniami
nad $\C$ takimi, że $\kappa(S)\ge0$ i $S,S'$
są biwymierne. Jeśli $K_S$ i $K_{S'}$ są nef,
to $S\cong S'$. Zatem każda klasa biwymierności
powierzchni o $\kappa\ge0$ ma, z dokładnością do
izomorfizmu, jeden gładki model kanonicznie minimalny.
\end{twierdzenie}
\begin{uwaga}[Dlaczego przypadek wymierny jest inny?]
Dla powierzchni o $\kappa=-\infty$ nie żądamy
kanonicznej nefności, bo $\PP^2$ ma $K=-3H$,
a wszystkie powierzchnie Hirzebrucha mają
$K\cdot f=-2$. Mogą jednak nie mieć krzywych
$(-1)$; otrzymane powierzchnie minimalne w sensie gładkim
nie muszą być izomorficzne. To wyjaśnia przykład
$\PP^2$ i $\PP^1\times\PP^1$.
\end{uwaga}

\section{Nefność, rozkład Zariskiego i granice teorii}
\label{sec:ga8-zariski-decomposition}

\begin{twierdzenie}[Rozkład Zariskiego na powierzchni — wynik zewnętrzny]
\label{tw:ga8-zariski-decomposition}
Niech $S$ będzie gładką rzutową powierzchnią nad
algebraicznie domkniętym ciałem charakterystyki zero,
a $D$ efektywnym dywizorem o współczynnikach wymiernych.
Istnieje jednoznaczny rozkład $D=P+N$, gdzie
$P,N$ są dywizorami o współczynnikach wymiernych,
$P$ jest nef, $N$ jest efektywny, a jeśli $N\ne0$,
to macierz przecięć składowych $N$ jest ujemnie
określona oraz $P\cdot N_i=0$ dla każdej składowej
nierozkładalnej $N_i$ dzielnika $N$.
\end{twierdzenie}

\begin{przyklad}[Rozkład Zariskiego obliczony do końca]
Niech $\pi:\widetilde S\to\PP^2$ będzie rozdmuchaniem
punktu, $H=\pi^*\mathcal O_{\PP^2}(1)$ i $E$
krzywą wyjątkową. Dla efektywnego dywizora
$D=H+E$ kładziemy $P=H$ i $N=E$.
Wówczas $H$ jest nef, $E$ jest efektywny,
$E^2=-1$ i $H\cdot E=0$. Spełnione są wszystkie
warunki twierdzenia, więc jest to jedyny rozkład.
Obserwujemy przy tym $D\cdot E=-1$,
choć dodatnia część $P$ nie ma ujemnych przecięć.
Nie oznacza to, że dodanie $E$ zwiększa liczbę
niezależnych przekrojów: część ujemna jest
właśnie stałą częścią odpowiednio wysokich
wielokrotności układu liniowego $|D|$.
\end{przyklad}

W \emph{wyższych wymiarach} nie wystarcza
kontraktowanie pojedynczych krzywych $(-1)$;
pojawiają się m.in. małe kontrakcje, przekształcenia
\emph{flip} oraz modele z osobliwościami.
Nie omawiamy ich tutaj. Naszym celem było
przeprowadzenie dwuwymiarowego algorytmu na
podstawie przecięć oraz pokazanie, które składniki
wymagają głębokich twierdzeń, a które wynikają
z elementarnego rachunku na powierzchni.

% Literatura do uzupelnienia MMP:
% J. Kollar, S. Mori, Birational Geometry of Algebraic Varieties;
% M. Reid, Chapters on Algebraic Surfaces;
% S. Frengley, Algebraic Surfaces (uniwersyteckie notatki online);
% Stacks Project, contractions and minimal models of surfaces.


\section{Podsumowanie geometrii biwymiernej i modeli minimalnych}
\label{sec:ga8-conclusion}

Przekształcenie Cremony dostarczyło kompletnego,
obliczalnego przykładu odwzorowania wymiernego,
które ma trzy punkty bazowe i kontraktuje trzy proste.
Rozdmuchanie tych punktów zamienia odwzorowanie
w morfizm, a graf przecięć sześciu krzywych
$(-1)$ pokazuje jego geometrię lokalną i globalną.
Wzory przecięć i klasy Cherna pozwalają śledzić
przy tym liczby geometryczne.

Zbudowaliśmy również język równoważności numerycznej,
stożka krzywych i nefności; otrzymaliśmy kryterium
krzywych $(-1)$ oraz dwuwymiarowy algorytm kontrakcji.
Wyraźnie wydzieliliśmy głębokie twierdzenia o stożku,
kontrakcjach i jednoznaczności modeli minimalnych od
rachunków wykonanych w skrypcie. W szczególności
nie utożsamiamy minimalności gładkiej powierzchni
z nefnością $K$ w przypadku $\kappa=-\infty$.

% Odniesienia kontrolne:
% J. Blanc, Introduction to Birational Geometry of Surfaces;
% I. Dolgachev, Classical Algebraic Geometry (Cremona transformations);
% Encyklopedia matematyki: Cremona transformations.

